\uuid{RaDb}
\exo7id{1589}
\titre{exo7 1589}
\auteur{ortiz}
\organisation{exo7}
\datecreate{1999-04-01}
\isIndication{false}
\isCorrection{false}
\chapitre{Réduction d'endomorphisme, polynôme annulateur}
\sousChapitre{Polynôme annulateur}
\cpge{mpsi-mp2i.calcul-matriciel-systemes.e, pcsi.determinants.c, ptsi.determinants.a, mpsi-mp2i.polynomes-fractions-rationnelles.a, pcsi.polynomes.a, ptsi.polynomes.a, mpsi-mp2i.changements-bases.c, mpsi-mp2i.determinants.d, pcsi.calcul-matriciel-systemes-lineaires.d, ptsi.calcul-matriciel-systemes-lineaires.d, mp-mpi.reduction-endomorphismes.g, mp-mpi.reduction-endomorphismes.i, mp-mpi.reduction-endomorphismes.c}
\module{Algèbre}
\niveau{L2}
\difficulte{}

\contenu{
\begin{enumerate}
    \item \question{Montrer que $P=(X-1)^2(X-2)$ est un polyn\^ome annulateur
de la matrice $A=\left(\begin{smallmatrix}
1&0&0\\0&1&0\\0&0&2\end{smallmatrix}\right)$ et en
d\'eduire le polyn\^ome minimal de la matrice $A$.}
    \item \question{Soit $B\in\mathcal{M}_2(\Cc)$. Calculer explicitement
$B^2-\text{tr}(B)\,B+\det(B)I_2$. En d\'eduire le
polyn\^ome minimal de la matrice
$B=\left(\begin{smallmatrix}
3&1\\-1&1\end{smallmatrix}\right)$.}
\end{enumerate}
}
